\documentclass[12pt]{article}
\usepackage[utf8]{inputenc}
\usepackage[T1]{fontenc}
\usepackage{hyperref}
\title{Tiwa, Northern Language Resources for K-12 Curriculum Development}
\author{Compiled by Claude Sonnet 5 (claude-sonnet-5) for Andr\'as Kornai}
\date{}
\begin{document}
\maketitle

\section{Introduction}
Northern Tiwa (ISO 639-3 twf; Glottocode nort1550) is a Kiowa-Tanoan language of New Mexico. Ethnologue treats Taos and Picuris as its two dialects but says they are not mutually intelligible. Other sources treat them as separate languages. Figures found in search: about 800 Taos speakers (Golla 2007) and about 225 Picuris speakers (Ethnologue 2007). UNESCO classes Picuris as Vulnerable. My search did not find a current dictionary or a named language programme for either variety. A summary of one result mentions a Taos Pueblo article on Tiwa language programmes for preschoolers, which I could not open.

\section{Primary Educational Resources}
None confirmed. The native-languages.org page below reportedly links to a Taos Pueblo article on preschool Tiwa programmes; I did not see its content.

\section{Linguistic and Reference Materials}
\begin{itemize}
\item Glottolog entry for Northern Tiwa, which lists grammatical and lexical sources: \url{https://glottolog.org/resource/languoid/id/nort1550}
\item WALS entry for Taos: \url{https://wals.info/languoid/lect/wals_code_twn}
\item Picuris language (Wikipedia): \url{https://en.wikipedia.org/wiki/Picuris_language}
\item Glottolog and WALS point to Trager's work on Taos grammar and Picuris phonology and to Zaharlick's 1977 study of Picuris syntax. I did not see those texts.
\end{itemize}

\section{Religious and Community Resources}
\begin{itemize}
\item Global Recordings Network, Taos audio Bible recordings: \url{https://www.globalrecordings.net/en/language/3788}
\end{itemize}

\section{Digital and Audio Resources}
\begin{itemize}
\item Endangered Languages Project entry: \url{https://endangeredlanguages.com/lang/6510}
\item Verbix page on Taos: \url{https://docs.verbix.com/EndangeredLanguages/Taos}
\item Verbix page on Picuris: \url{https://docs.verbix.com/EndangeredLanguages/Picuris}
\item Native-languages.org Tiwa overview: \url{https://www.native-languages.org/tiwa.htm}
\end{itemize}

\section{Curriculum Development Context}
Adults still use the language as a first language, but I found no school materials and no confirmed pueblo-run programme. Taos and Picuris would need separate materials, and any work must be led by the two pueblos, who control their languages. The next step is to ask the Taos Pueblo and Picuris Pueblo education offices what exists. I have no contact details to give.
\end{document}